\section{Metric TSP and Eulerian Structure}\label{tsp:main}

\subsection{General TSP, metric TSP, and the shortcut rule}
The traveling salesperson problem asks for a minimum-cost Hamiltonian
cycle: visit every vertex exactly once and return to the start.  Throughout
the approximation algorithms below, the input is a complete undirected
graph with nonnegative symmetric costs satisfying the triangle inequality
\[
c(u,w)\le c(u,v)+c(v,w).
\]
These are \emph{metric} costs.  Completeness supplies each shortcut edge;
triangle inequality ensures that shortcutting a walk never increases its
cost.  Repeatedly bypass already visited vertices, retaining the final
return to the start, to turn a spanning closed walk into a Hamiltonian cycle.

\paragraph{Why the metric assumption matters.}
General TSP has no polynomial-time approximation with any fixed finite
factor $\rho\ge1$, unless $\PTIME=\NP$.  Reduce Hamiltonian Cycle on an
unweighted graph $G$ with $n$ vertices: complete the graph, give original
edges cost $1$ and nonedges cost an integer $B>\rho n$.  A YES instance has
$\OPT=n$, so a $\rho$-approximation returns cost at most $\rho n$.  In a NO
instance every tour uses a nonedge, hence has cost at least $B>\rho n$.
The algorithm would distinguish the two cases in polynomial time.  For
fixed $\rho$, $B$ has polynomial encoding length.  These constructed costs
need not satisfy triangle inequality; a cost-$B$ shortcut can replace a
two-edge detour of cost $2$.

Metric TSP remains NP-hard: assign costs $1$ to edges of a Hamiltonian Cycle
instance and $2$ to nonedges.  Every triangle satisfies the metric inequality,
and a tour of cost $n$ exists exactly when the original graph is Hamiltonian.
This proves exact hardness, not the absence of constant-factor approximation.

\subsection{Euler circuits, open trails, and orientation}\label{tsp:euler}
An Euler trail uses every \emph{edge occurrence} exactly once; vertices may
repeat.  Parallel edges are distinct occurrences.  An Euler circuit is a
closed Euler trail.  A Hamiltonian cycle instead visits every vertex once.

\begin{theorem}[Euler's criterion]
An undirected multigraph with at least one edge has an Euler circuit iff
all its nonisolated vertices are connected and every vertex has even degree.
It has an Euler trail between distinct vertices $s,t$ iff all nonisolated
vertices are connected and precisely $s,t$ have odd degree.
\end{theorem}
\begin{proof}
In a closed traversal, each arrival at a vertex is paired with a departure,
so every degree is even.  For an open traversal only the start and end have
one unpaired incident edge.

Conversely, start at any nonisolated vertex and keep following unused edges.
An even-degree vertex other than the start cannot be the first place where
this trail gets stuck: after arriving, an unused departure remains.  Thus
the trail closes.  If edges remain, connectivity gives a vertex of the
current circuit incident to unused edges.  The unused-edge degrees are
still even, so construct another closed trail there and splice it into the
first.  Repeat until every edge is used (Hierholzer's algorithm).
For the open case add one extra $\{s,t\}$ edge, find an Euler circuit,
rotate it to place that edge last, and delete the extra edge.
\end{proof}
Using adjacency lists and marking each edge occurrence, Hierholzer's
algorithm runs in $O(|V|+|E|)$ time.  The handshaking identity
$\sum_v\deg(v)=2|E|$ also shows that the number of odd-degree vertices is
even.  Adding one edge toggles the parity of its two endpoints.

\paragraph{Orienting an even-degree connected graph.}
Find an Euler circuit and orient each edge in the direction in which the
circuit traverses it.  The resulting directed graph is strongly connected:
from any vertex $u$, follow the directed circuit cyclically until any desired
vertex $v$ is reached.  Also $\deg^+(v)=\deg^-(v)$ at every vertex.
This proof concerns the nonisolated vertices; a disconnected even-degree
graph does not become strongly connected by orienting its components.

\subsection{Double-tree: a 2-approximation}\label{tsp:doubletree}
\begin{enumerate}
  \item Compute a minimum spanning tree $T$.
  \item Duplicate every edge of $T$, retaining both copies in a multigraph.
  \item Find an Euler circuit in this connected even-degree multigraph.
  \item Shortcut repeated vertex visits to obtain the output tour $H$.
\end{enumerate}
\textbf{Feasibility.} Doubling makes every degree even; the spanning tree
ensures connectivity.  The Euler circuit visits all vertices, and
shortcutting yields a Hamiltonian cycle in the complete graph.

\textbf{Lower bound and guarantee.} Deleting any edge of an optimal tour
leaves a spanning tree, so $c(T)\le\OPT$.  Therefore
\[
\boxed{\ c(H)\le 2c(T)\le2\OPT.\ }
\]
The first inequality uses triangle inequality; the second uses the MST lower
bound.  Dense Prim takes $O(n^2)$ time on the complete graph; doubling,
Euler traversal, and shortcutting take $O(n)$ further time.

\paragraph{Four-vertex metric counterexample to optimality.}
Use the shortest-path metric of the unit-cost cycle $A,B,D,C,A$:
\[
\begin{array}{c|rrrr}
c & A&B&C&D\\\hline
A&0&1&1&2\\
B&1&0&2&1\\
C&1&2&0&1\\
D&2&1&1&0
\end{array}
\]
The tree $\{AB,AC,CD\}$ is an MST of cost $3$.  The doubled tree admits
Euler circuit $A,B,A,C,D,C,A$.  First-visit shortcutting gives
$A,B,C,D,A$, of cost $1+2+1+2=6$.  The original cycle has cost $4$ and is
optimal because every tour has four edges, each of cost at least $1$.
Thus a valid double-tree execution need not be optimal, even on a metric.

\subsection{Christofides: a \texorpdfstring{$3/2$}{3/2}-approximation}\label{tsp:christofides}
\begin{enumerate}
  \item Compute an MST $T$, and let $O=\{v:\deg_T(v)\text{ is odd}\}$.
  \item Find a \emph{minimum-weight perfect matching} $M$ on the complete
        graph induced by $O$, using the original metric costs.
  \item Form the multigraph $T\cup M$, find an Euler circuit, and shortcut
        repeated vertices to produce a tour $H$.
\end{enumerate}
\textbf{Feasibility.} The set $O$ has even cardinality.  Every vertex in
$O$ gains exactly one incident matching edge and becomes even; all other
degrees stay even.  The tree ensures connectivity.  Keep parallel copies
if a matching edge is also in $T$.

\begin{lemma}[Matching budget]
The matching chosen above has cost $c(M)\le\OPT/2$.
\end{lemma}
\begin{proof}
Take an optimal tour and shortcut it to retain only the vertices of $O$.
The resulting even-length cyclic walk has cost at most $\OPT$.  Its
alternating edges split into two perfect matchings $M_1,M_2$ on $O$.
Therefore, by minimum weight of $M$,
\[
c(M)\le\min\{c(M_1),c(M_2)\}
\le\frac{c(M_1)+c(M_2)}2\le\frac{\OPT}2.
\]
If $O$ has two vertices, the restricted cyclic walk traverses their edge
twice; its two edge occurrences give the same one-edge matching.  If
$O=\varnothing$, use the empty matching.
\end{proof}
Combining the spanning and matching budgets with metric shortcutting gives
\[
\boxed{\ c(H)\le c(T)+c(M)\le\OPT+\tfrac12\OPT
       =\tfrac32\OPT.\ }
\]
All steps are polynomial time; minimum-weight perfect matching is the main
additional subroutine beyond double-tree.  A maximal matching is insufficient:
it may leave odd vertices unmatched, and an arbitrary perfect matching has
no $\OPT/2$ cost guarantee.

\subsection{Held--Karp exact dynamic programming}\label{tsp:heldkarp}
This exact algorithm works for arbitrary tour costs; triangle inequality
is unnecessary.  Fix a root $r$.  For $r\in S\subseteq V$ and $v\in S$,
let $D[S,v]$ be the minimum cost of a path starting at $r$, ending at $v$,
and visiting exactly the vertices in $S$, once each.

\textbf{Bases.} $D[\{r\},r]=0$.  Set all impossible states to $\infty$,
including $D[S,r]$ for $|S|>1$ and states with $r\notin S$ or $v\notin S$.
For $v\ne r$,
\[
\boxed{\ D[S,v]=\min_{u\in S\setminus\{v\}}
        \bigl(D[S\setminus\{v\},u]+c(u,v)\bigr).\ }
\]
\textbf{Why the recurrence is exact.} The penultimate vertex $u$ partitions
all feasible paths.  Removing their last edge leaves precisely a feasible
path for the smaller state; conversely, appending $(u,v)$ to any such path
gives a feasible path for $(S,v)$.

Evaluate states in increasing $|S|$.  For $n\ge2$, close the path by returning
to $r$:
\[
\OPT=\min_{v\ne r}\bigl(D[V,v]+c(v,r)\bigr).
\]
Store a minimizing predecessor with each state; choose the minimizing final
vertex, backtrack while removing the last vertex from $S$, reverse the
resulting sequence, and append $r$.  There are $O(n2^n)$ states with $O(n)$
transitions each, giving $O(n^2 2^n)$ time and $O(n2^n)$ space, including
predecessors.  This is exponential, although faster than enumerating tours.
